%=============================================================================
% Chapter 4 Exercises
%=============================================================================
\newpage
\section{Chapter 4 Exercises [1060]}

%-----------------------------------------------------------------------------
\subsection{[20] Control Signals for \texttt{lw2}}
\hypertarget{prob-4.1}{}
\noindent$\langle$\S4.3$\rangle$\medskip

\noindent Consider the following new instruction:
\begin{verbatim}
    lw2 $t1, $s1, $s2    # $t1 = Mem[$s1 + Mem[$s2]]
\end{verbatim}

\begin{enumerate}
  \phantomsection\addcontentsline{toc}{subsubsection}{4.1.1~[5]}
  \item[4.1.1] [5] What are the values of control signals generated by the
        control in Figure 4.10 for this instruction?
  \phantomsection\addcontentsline{toc}{subsubsection}{4.1.2~[5]}
  \item[4.1.2] [5] Which resources (blocks) perform a useful function for
        this instruction?
  \phantomsection\addcontentsline{toc}{subsubsection}{4.1.3~[10]}
  \item[4.1.3] [10] Which resources (blocks) produce no output for this
        instruction? Which resources produce output that is not used?
\end{enumerate}

\problembreak

%-----------------------------------------------------------------------------
\subsection{[10] Don't Cares in Control}
\hypertarget{prob-4.2}{}
\noindent$\langle$\S4.4$\rangle$\medskip

\noindent Explain each of the ``don't cares'' in Figure 4.18.

\problembreak

%-----------------------------------------------------------------------------
\subsection{[15] Instruction Mix Resource Usage}
\hypertarget{prob-4.3}{}
\noindent$\langle$\S4.4$\rangle$\medskip

\noindent Consider the following instruction mix:

\begin{center}
\begin{tabular}{cccccc}
\toprule
\textbf{R-type} & \textbf{I-type (non-lw)} & \textbf{Load} &
\textbf{Store} & \textbf{Branch} & \textbf{Jump} \\
\midrule
24\% & 28\% & 25\% & 10\% & 11\% & 2\% \\
\bottomrule
\end{tabular}
\end{center}

\begin{enumerate}
  \phantomsection\addcontentsline{toc}{subsubsection}{4.3.1~[5]}
  \item[4.3.1] [5] What fraction of all instructions use data memory?
  \phantomsection\addcontentsline{toc}{subsubsection}{4.3.2~[5]}
  \item[4.3.2] [5] What fraction of all instructions use the sign extend?
  \phantomsection\addcontentsline{toc}{subsubsection}{4.3.3~[5]}
  \item[4.3.3] [5] What is the sign extend doing during cycles in which its
        output is not needed?
\end{enumerate}

\problembreak

%-----------------------------------------------------------------------------
\subsection{[10] Stuck-at-0 Fault Analysis}
\hypertarget{prob-4.4}{}
\noindent$\langle$\S4.4$\rangle$\medskip

\noindent When silicon chips are fabricated, defects in materials (e.g.,
silicon) and manufacturing errors can result in defective circuits. A very
common defect is for one signal wire to get ``broken'' and always register a
logical 0. This is often called a ``stuck-at-0'' fault.

\begin{enumerate}
  \phantomsection\addcontentsline{toc}{subsubsection}{4.4.1~[5]}
  \item[4.4.1] [5] Which instruction(s) fail to operate correctly if the
        \texttt{RegDst} wire is stuck at 0?
  \phantomsection\addcontentsline{toc}{subsubsection}{4.4.2~[5]}
  \item[4.4.2] [5] Which instruction(s) fail to operate correctly if the
        \texttt{ALUSrc} wire is stuck at 0?
\end{enumerate}

\problembreak

%-----------------------------------------------------------------------------
\subsection{[45] Trace Instruction Through Datapath}
\hypertarget{prob-4.5}{}
\noindent$\langle$\S4.4$\rangle$\medskip

\noindent Consider the single-cycle datapath in Figure 4.17. Trace the
execution of the following instruction through the datapath:
\begin{verbatim}
    add $s0, $s1, $s2
\end{verbatim}

\begin{enumerate}
  \phantomsection\addcontentsline{toc}{subsubsection}{4.5.1~[10]}
  \item[4.5.1] [10] What are the values of control signals generated by the
        control unit?
  \phantomsection\addcontentsline{toc}{subsubsection}{4.5.2~[5]}
  \item[4.5.2] [5] What is the new PC address after this instruction is
        executed? Highlight the path through which this value is determined.
  \phantomsection\addcontentsline{toc}{subsubsection}{4.5.3~[10]}
  \item[4.5.3] [10] For each mux, show the values of its inputs and outputs
        during the execution of this instruction. List the values of all
        register outputs.
  \phantomsection\addcontentsline{toc}{subsubsection}{4.5.4~[10]}
  \item[4.5.4] [10] What are the input values for the ALU and the two
        read data outputs?
  \phantomsection\addcontentsline{toc}{subsubsection}{4.5.5~[10]}
  \item[4.5.5] [10] What are the values of all inputs for the registers?
\end{enumerate}

\problembreak

%-----------------------------------------------------------------------------
\subsection{[15] Adding I-type Instructions}
\hypertarget{prob-4.6}{}
\noindent$\langle$\S4.4$\rangle$\medskip

\noindent Section 4.1 does not discuss I-type instructions that are not
load/store.

\begin{enumerate}
  \phantomsection\addcontentsline{toc}{subsubsection}{4.6.1~[5]}
  \item[4.6.1] [5] What additional logic blocks, if any, are needed to add
        I-type instructions to the CPU shown in Figure 4.17? Add any
        necessary logic blocks to Figure 4.12 and explain their purpose.
  \phantomsection\addcontentsline{toc}{subsubsection}{4.6.2~[10]}
  \item[4.6.2] [10] List the values of the signals generated by the control
        unit for each I-type instruction. Explain the reasoning for any
        ``Don't care'' control signals.
\end{enumerate}

\problembreak

%-----------------------------------------------------------------------------
\subsection{[30] Instruction Latencies}
\hypertarget{prob-4.7}{}
\noindent$\langle$\S4.4$\rangle$\medskip

\noindent Problems in this exercise assume that the logic blocks used to
implement a processor's datapath have the following latencies:

\begin{center}
\begin{tabular}{ccccccccc}
\toprule
\textbf{I-Mem} & \textbf{Mux} & \textbf{ALU} & \textbf{Adder} &
\textbf{Gate} & \textbf{Reg} & \textbf{Reg File} &
\textbf{Sign ext.} & \textbf{Ctrl} \\
\midrule
200ps & 30ps & 120ps & 150ps & 5ps & 90ps & 150ps & 10ps & 50ps \\
\bottomrule
\end{tabular}
\end{center}

\noindent ``Register read'' is the time after the rising clock edge for the
new register value to appear on the output. ``Register setup'' is the time a
register's data input must be stable before the rising edge of the clock.

\begin{enumerate}
  \phantomsection\addcontentsline{toc}{subsubsection}{4.7.1~[5]}
  \item[4.7.1] [5] What is the latency of an R-type instruction (i.e., how
        long must the clock period be to ensure that this instruction works
        correctly)?
  \phantomsection\addcontentsline{toc}{subsubsection}{4.7.2~[10]}
  \item[4.7.2] [10] What is the latency of a \texttt{lw} instruction?
  \phantomsection\addcontentsline{toc}{subsubsection}{4.7.3~[10]}
  \item[4.7.3] [10] What is the latency of an arithmetic, logical, or shift
        I-type (non-load) instruction?
  \phantomsection\addcontentsline{toc}{subsubsection}{4.7.4~[5]}
  \item[4.7.4] [5] What is the minimum clock period for this CPU?
\end{enumerate}

\problembreak

%-----------------------------------------------------------------------------
\subsection{[10] Variable Clock Cycle Speedup}
\hypertarget{prob-4.8}{}
\noindent$\langle$\S4.4$\rangle$\medskip

\noindent Suppose you could build a CPU where the clock cycle time was
different for each instruction. What would be the speedup of this CPU over
the CPU presented in Figure 4.23 given the instruction mix from Exercise 4.3?

\problembreak

%-----------------------------------------------------------------------------
\subsection{[25] Multiplier Addition to CPU}
\hypertarget{prob-4.9}{}
\noindent$\langle$\S4.4$\rangle$\medskip

\noindent Consider the addition of a multiplier to the CPU shown in Figure
4.21. This addition will add 300 ps to the latency of the ALU, but will
reduce the number of instructions by 5\% (because there will no longer be a
need to simulate the multiply instruction).

\begin{enumerate}
  \phantomsection\addcontentsline{toc}{subsubsection}{4.9.1~[5]}
  \item[4.9.1] [5] What is the clock cycle time with and without this
        improvement?
  \phantomsection\addcontentsline{toc}{subsubsection}{4.9.2~[10]}
  \item[4.9.2] [10] What is the speedup achieved by adding this improvement?
  \phantomsection\addcontentsline{toc}{subsubsection}{4.9.3~[10]}
  \item[4.9.3] [10] What is the slowest the new ALU can be and still result
        in improved performance?
\end{enumerate}

\problembreak

%-----------------------------------------------------------------------------
\subsection{[25] Register File Expansion Cost}
\hypertarget{prob-4.10}{}
\noindent$\langle$\S4.4$\rangle$\medskip

\noindent When processor designers consider a possible improvement, the
decision usually depends on the cost/performance tradeoff. Assume we are
beginning with the datapath from Figure 4.21 and the latencies from Exercise
4.7. Suppose doubling the number of general purpose registers from 32 to 64
would reduce the number of \texttt{lw} and \texttt{sw} instructions executed
by 12\%, but increase the latency of the register file from 150 ps to 160 ps
and double the cost from 200 to 400.

\begin{enumerate}
  \phantomsection\addcontentsline{toc}{subsubsection}{4.10.1~[5]}
  \item[4.10.1] [5] What is the speedup achieved by adding this improvement?
  \phantomsection\addcontentsline{toc}{subsubsection}{4.10.2~[10]}
  \item[4.10.2] [10] Compare the change in performance to the change in cost.
  \phantomsection\addcontentsline{toc}{subsubsection}{4.10.3~[10]}
  \item[4.10.3] [10] Describe a situation where it makes sense to add more
        registers and a situation where it doesn't.
\end{enumerate}

\problembreak

%-----------------------------------------------------------------------------
\subsection{[20] Load With Increment Instruction}
\hypertarget{prob-4.11}{}
\noindent$\langle$\S4.4$\rangle$\medskip

\noindent Examine the difficulty of adding a proposed ``Load With Increment''
instruction to MIPS:
\begin{verbatim}
    lwi $t1, 0($s1)    # $t1 = Mem[$s1]; $s1 = $s1 + 4
\end{verbatim}

\begin{enumerate}
  \phantomsection\addcontentsline{toc}{subsubsection}{4.11.1~[5]}
  \item[4.11.1] [5] Which new functional blocks (if any) do we need?
  \phantomsection\addcontentsline{toc}{subsubsection}{4.11.2~[5]}
  \item[4.11.2] [5] Which existing functional blocks (if any) require
        modification?
  \phantomsection\addcontentsline{toc}{subsubsection}{4.11.3~[5]}
  \item[4.11.3] [5] What new data paths do we need (if any)?
  \phantomsection\addcontentsline{toc}{subsubsection}{4.11.4~[5]}
  \item[4.11.4] [5] What new signals do we need (if any) from the control
        unit to support this instruction?
\end{enumerate}

\problembreak

%-----------------------------------------------------------------------------
\subsection{[30] Swap Instruction}
\hypertarget{prob-4.12}{}
\noindent$\langle$\S4.4$\rangle$\medskip

\noindent Examine the difficulty of adding a proposed \texttt{swap}
instruction to MIPS:
\begin{verbatim}
    swap $s1, $s2       # temp = $s1; $s1 = $s2; $s2 = temp
\end{verbatim}

\begin{enumerate}
  \phantomsection\addcontentsline{toc}{subsubsection}{4.12.1~[5]}
  \item[4.12.1] [5] Which new functional blocks (if any) do we need?
  \phantomsection\addcontentsline{toc}{subsubsection}{4.12.2~[10]}
  \item[4.12.2] [10] Which existing functional blocks (if any) require
        modification?
  \phantomsection\addcontentsline{toc}{subsubsection}{4.12.3~[5]}
  \item[4.12.3] [5] What new data paths do we need (if any)?
  \phantomsection\addcontentsline{toc}{subsubsection}{4.12.4~[5]}
  \item[4.12.4] [5] What new signals do we need (if any) from the control
        unit?
  \phantomsection\addcontentsline{toc}{subsubsection}{4.12.5~[5]}
  \item[4.12.5] [5] Modify Figure 4.21 to demonstrate an implementation of
        this new instruction.
\end{enumerate}

\problembreak

%-----------------------------------------------------------------------------
\subsection{[35] Store Sum Instruction}
\hypertarget{prob-4.13}{}
\noindent$\langle$\S4.4$\rangle$\medskip

\noindent Examine the difficulty of adding a proposed ``Store Sum''
instruction to MIPS:
\begin{verbatim}
    ss $s1, $s2, 100($s3)  # Mem[$s3+100] = $s1 + $s2
\end{verbatim}

\begin{enumerate}
  \phantomsection\addcontentsline{toc}{subsubsection}{4.13.1~[10]}
  \item[4.13.1] [10] Which new functional blocks (if any) do we need?
  \phantomsection\addcontentsline{toc}{subsubsection}{4.13.2~[10]}
  \item[4.13.2] [10] Which existing functional blocks (if any) require
        modification?
  \phantomsection\addcontentsline{toc}{subsubsection}{4.13.3~[5]}
  \item[4.13.3] [5] What new data paths do we need (if any)?
  \phantomsection\addcontentsline{toc}{subsubsection}{4.13.4~[5]}
  \item[4.13.4] [5] What new signals do we need (if any) from the control
        unit?
  \phantomsection\addcontentsline{toc}{subsubsection}{4.13.5~[5]}
  \item[4.13.5] [5] Modify Figure 4.21 to demonstrate an implementation of
        this new instruction.
\end{enumerate}

\problembreak

%-----------------------------------------------------------------------------
\subsection{[5] Imm Gen Block Critical Path}
\hypertarget{prob-4.14}{}
\noindent$\langle$\S4.5$\rangle$\medskip

\noindent For which instructions (if any) is the Imm Gen block on the
critical path?

\problembreak

%-----------------------------------------------------------------------------
\subsection{[25] No-Offset Load/Store Redesign}
\hypertarget{prob-4.15}{}
\noindent$\langle$\S4.4$\rangle$\medskip

\noindent The \texttt{lw} instruction is the instruction with the longest
latency on the CPU from Section 4.4. If we modified \texttt{lw} and
\texttt{sw} so that they have no offset (i.e., the address to be loaded
from/stored to must be calculated and placed in \texttt{rs} before calling
\texttt{lw}/\texttt{sw}), then no instruction would use both the ALU and data
memory. This would allow us to reduce the clock cycle time. However, it would
also increase the number of instructions, because many \texttt{lw} and
\texttt{sw} instructions would need to be replaced with \texttt{add} or
combined instructions.

\begin{enumerate}
  \phantomsection\addcontentsline{toc}{subsubsection}{4.15.1~[5]}
  \item[4.15.1] [5] What would the new clock cycle time be?
  \phantomsection\addcontentsline{toc}{subsubsection}{4.15.2~[10]}
  \item[4.15.2] [10] Would a program with the instruction mix presented in
        Exercise 4.7 run faster or slower on this new CPU? By how much? (For
        simplicity, assume any \texttt{lw} or \texttt{sw} is replaced with a
        sequence of two instructions.)
  \phantomsection\addcontentsline{toc}{subsubsection}{4.15.3~[5]}
  \item[4.15.3] [5] What is the primary factor that influences whether a
        program will run faster or slower on the new CPU?
  \phantomsection\addcontentsline{toc}{subsubsection}{4.15.4~[5]}
  \item[4.15.4] [5] Do you consider the original CPU (as shown in Figure
        4.21) a better overall design, or do you consider the new CPU a
        better overall design? Why?
\end{enumerate}

\problembreak

%-----------------------------------------------------------------------------
\subsection{[45] Pipelining and Clock Cycle Time}
\hypertarget{prob-4.16}{}
\noindent$\langle$\S4.6$\rangle$\medskip

\noindent In this exercise, we examine how pipelining affects the clock cycle
time of the processor. Assume that individual stages of the datapath have
the following latencies:

\begin{center}
\begin{tabular}{ccccc}
\toprule
\textbf{IF} & \textbf{ID} & \textbf{EX} & \textbf{MEM} & \textbf{WB} \\
\midrule
250ps & 350ps & 150ps & 300ps & 200ps \\
\bottomrule
\end{tabular}
\end{center}

\noindent Also, assume that instructions executed by the processor are broken
down as follows: ALU/Logic 40\%, Jump/Branch 20\%, Load 20\%, Store 20\%.

\begin{enumerate}
  \phantomsection\addcontentsline{toc}{subsubsection}{4.16.1~[5]}
  \item[4.16.1] [5] What is the clock cycle time in a pipelined and
        non-pipelined processor?
  \phantomsection\addcontentsline{toc}{subsubsection}{4.16.2~[10]}
  \item[4.16.2] [10] What is the total latency of an \texttt{lw} instruction
        in a pipelined and non-pipelined processor?
  \phantomsection\addcontentsline{toc}{subsubsection}{4.16.3~[10]}
  \item[4.16.3] [10] If we can split one stage of the pipelined datapath
        into two new stages, each with half the latency of the original
        stage, which stage would you split and what is the new clock cycle
        time of the processor?
  \phantomsection\addcontentsline{toc}{subsubsection}{4.16.4~[10]}
  \item[4.16.4] [10] Assuming there are no stalls or hazards, what is the
        utilization of the data memory?
  \phantomsection\addcontentsline{toc}{subsubsection}{4.16.5~[10]}
  \item[4.16.5] [10] Assuming there are no stalls or hazards, what is the
        utilization of the write register port of the ``Registers'' unit?
\end{enumerate}

\problembreak

%-----------------------------------------------------------------------------
\subsection{[10] Minimum Cycles for N Instructions}
\hypertarget{prob-4.17}{}
\noindent$\langle$\S4.6$\rangle$\medskip

\noindent What is the minimum number of cycles needed to completely execute a
program of $N$ instructions on a CPU with a 4-stage pipeline? Justify your
answer.

\problembreak

%-----------------------------------------------------------------------------
\subsection{[5] Pipeline without Hazard Detection}
\hypertarget{prob-4.18}{}
\noindent$\langle$\S4.6$\rangle$\medskip

\noindent Assume that \texttt{\$s0} is initialized to 11 and \texttt{\$s1} is
initialized to 22. Suppose you executed the code below on a version of the
pipeline from Section 4.6 that does not handle data hazards (i.e., the
programmer is responsible for addressing data hazards by inserting
\texttt{nop} instructions where necessary). What would the final values of
registers be?
\begin{verbatim}
    add $s0, $s0, $s1
    add $s2, $s0, $s1
\end{verbatim}

\problembreak

%-----------------------------------------------------------------------------
\subsection{[10] Pipeline with Same-Cycle Forwarding}
\hypertarget{prob-4.19}{}
\noindent$\langle$\S4.6$\rangle$\medskip

\noindent Assume that \texttt{\$s0} is initialized to 11 and \texttt{\$s1} is
initialized to 22. Suppose you executed the code below on a version of the
pipeline from Section 4.6 that does not handle data hazards. See
Section 4.8: an ID stage will return the results of a WB state occurring
during the same cycle. What would the final values of registers be?
\begin{verbatim}
    add $s0, $s0, $s1
    nop
    nop
    add $s2, $s0, $s1
\end{verbatim}

\problembreak

%-----------------------------------------------------------------------------
\subsection{[5] Adding NOPs for Data Hazards}
\hypertarget{prob-4.20}{}
\noindent$\langle$\S4.6$\rangle$\medskip

\noindent Consider a version of the pipeline that does not handle data
hazards. Add \texttt{nop} instructions to the code below so that it will run
correctly on a pipeline that does not handle data hazards:
\begin{verbatim}
    add $s0, $s1, $s5
    add $s2, $s0, $s5
    add $s4, $s0, $s1
\end{verbatim}

\problembreak

%-----------------------------------------------------------------------------
\subsection{[5] Forwarding Hardware Speedup}
\hypertarget{prob-4.21}{}
\noindent$\langle$\S4.6$\rangle$\medskip

\noindent Suppose that the clock cycle time of the pipeline without
forwarding is 250 ps. Suppose also that adding forwarding hardware will
reduce the number of \texttt{nop}s from 5\% to 2\%, but increase the cycle
time to 300 ps. What is the speedup of the new pipeline compared to the one
without forwarding?

\problembreak

%-----------------------------------------------------------------------------
\subsection{[20] Structural Hazard: Unified Memory}
\hypertarget{prob-4.22}{}
\noindent$\langle$\S4.6$\rangle$\medskip

\noindent Consider the fragment of MIPS assembly below:
\begin{verbatim}
    lw  $s5, 8($s3)
    sw  $s4, 0($s3)
    add $s2, $s4, $s1
\end{verbatim}

\noindent Suppose we modify the pipeline so that it has only one memory (that
handles both instructions and data). In this case, there will be a structural
hazard every time a program needs to fetch an instruction during the same
cycle in which another instruction accesses data.

\begin{enumerate}
  \phantomsection\addcontentsline{toc}{subsubsection}{4.22.1~[5]}
  \item[4.22.1] [5] Draw a pipeline diagram to show where the code above
        stalls.
  \phantomsection\addcontentsline{toc}{subsubsection}{4.22.2~[5]}
  \item[4.22.2] [5] In general, is it possible to reduce the number of
        stalls/nops resulting from this structural hazard by reordering code?
  \phantomsection\addcontentsline{toc}{subsubsection}{4.22.3~[5]}
  \item[4.22.3] [5] Must this structural hazard be handled in hardware? Can
        you eliminate it by adding \texttt{nop}s to the code? If so, explain
        how. If not, explain why not.
  \phantomsection\addcontentsline{toc}{subsubsection}{4.22.4~[5]}
  \item[4.22.4] [5] Approximately how many stalls would you expect this
        structural hazard to generate in a typical program? (Use the
        instruction mix from Exercise 4.3.)
\end{enumerate}

\problembreak

%-----------------------------------------------------------------------------
\subsection{[20] Four-Stage Pipeline}
\hypertarget{prob-4.23}{}
\noindent$\langle$\S4.6$\rangle$\medskip

\noindent If we change the load/store instructions to use a register (without
an offset) as the address, those instructions no longer need to use the ALU.
(See Exercise 4.15.) As a result, the MEM and EX stages can be overlapped and
the pipeline has only four stages.

\begin{enumerate}
  \phantomsection\addcontentsline{toc}{subsubsection}{4.23.1~[10]}
  \item[4.23.1] [10] How will the reduction in pipeline depth affect
        performance?
  \phantomsection\addcontentsline{toc}{subsubsection}{4.23.2~[5]}
  \item[4.23.2] [5] How might this change improve the pipeline performance?
  \phantomsection\addcontentsline{toc}{subsubsection}{4.23.3~[5]}
  \item[4.23.3] [5] How might this change degrade the performance?
\end{enumerate}

\problembreak

%-----------------------------------------------------------------------------
\subsection{[10] Pipeline Diagram Selection}
\hypertarget{prob-4.24}{}
\noindent$\langle$\S4.6$\rangle$\medskip

\noindent Which of the two pipeline diagrams below better represents the
operation of the pipeline's hazard detection unit? Why?

\begin{enumerate}
  \item[a.]
\begin{verbatim}
  ld  x1, x12:    IF ID EX ME WB
  add x5, x1, x17:   IF ID EX ME WB
\end{verbatim}

  \item[b.]
\begin{verbatim}
  ld  x1, x12:    IF ID EX ME WB
  add x5, x1, x17:      IF ID EX ME WB
\end{verbatim}
\end{enumerate}

\problembreak

%-----------------------------------------------------------------------------
\subsection{[10] Loop Pipeline Execution}
\hypertarget{prob-4.25}{}
\noindent$\langle$\S4.6$\rangle$\medskip

\noindent Consider the following loop:
\begin{verbatim}
    LOOP: lw   $s0, 0($s3)
          add  $s1, $s0, $s2
          sw   $s1, 0($s3)
          addi $s3, $s3, -4
          bne  $s3, $zero, LOOP
\end{verbatim}

\noindent Assume that perfect branch prediction is used (no stalls due to
control hazards), that there are no delay slots, that the pipeline has full
forwarding support, and that branches are resolved in the EX stage.

\begin{enumerate}
  \phantomsection\addcontentsline{toc}{subsubsection}{4.25.1~[10]}
  \item[4.25.1] [10] Show a pipeline execution diagram for the first two
        iterations of this loop.
\end{enumerate}

\problembreak

%-----------------------------------------------------------------------------
\subsection{[50] RAW Dependencies and Forwarding}
\hypertarget{prob-4.26}{}
\noindent$\langle$\S4.8$\rangle$\medskip

\noindent This exercise refers to pipelined datapath from Figure 4.33. The
following fractions of instructions have particular types of RAW data
dependence (identified by the stage that produces the result and the
instruction that consumes it):

\begin{center}
\begin{tabular}{cccccc}
\toprule
\textbf{EX to 1st} & \textbf{MEM to 1st} & \textbf{EX to 2nd} &
\textbf{MEM to 2nd} & \textbf{EX to 1st \&} & \textbf{MEM to 1st \&} \\
& & & & \textbf{EX to 2nd} & \textbf{MEM to 2nd} \\
\midrule
20\% & 5\% & 20\% & 10\% & 10\% & 5\% \\
\bottomrule
\end{tabular}
\end{center}

\begin{enumerate}
  \phantomsection\addcontentsline{toc}{subsubsection}{4.26.1~[5]}
  \item[4.26.1] [5] For each RAW dependency listed above, give a sequence of
        at least three assembly statements that exhibits that dependence.
  \phantomsection\addcontentsline{toc}{subsubsection}{4.26.2~[5]}
  \item[4.26.2] [5] For each RAW dependency above, how many \texttt{nop}(s)
        would need to be inserted to allow your code from 4.26.1 to run
        correctly on a pipeline with no forwarding or hazard detection?
  \phantomsection\addcontentsline{toc}{subsubsection}{4.26.3~[10]}
  \item[4.26.3] [10] Write a sequence of three assembly instructions so that,
        when you consider each instruction independently, the sum of the
        stalls is larger than the number of stalls the sequence actually
        needs to avoid data hazards.
  \phantomsection\addcontentsline{toc}{subsubsection}{4.26.4~[5]}
  \item[4.26.4] [5] Assuming no other hazards, what is the CPI of the
        program if we use only EX-EX forwarding or only MEM-EX forwarding?
  \phantomsection\addcontentsline{toc}{subsubsection}{4.26.5~[5]}
  \item[4.26.5] [5] What is the CPI if we use full forwarding (forward from
        both EX/MEM and MEM/WB)? What percent of cycles is stalled?
  \phantomsection\addcontentsline{toc}{subsubsection}{4.26.6~[10]}
  \item[4.26.6] [10] What is the speedup achieved by each type of forwarding
        (EX/MEM, MEM/WB, full) compared to a pipeline with no forwarding?
  \phantomsection\addcontentsline{toc}{subsubsection}{4.26.7~[5]}
  \item[4.26.7] [5] What would be the additional speedup (relative to the
        fastest processor from 4.26.6) if we added ``time-travel''
        forwarding that eliminates all data hazards? Assume that the
        yet-to-be-invented time-travel circuitry adds 100 ps to the latency
        of the full-forwarding EX stage.
  \phantomsection\addcontentsline{toc}{subsubsection}{4.26.8~[5]}
  \item[4.26.8] [5] The table of hazard types has separate entries for ``EX
        to 1st'' and ``EX to 1st and EX to 2nd''. Why is there no entry for
        ``MEM to 1st'' and ``MEM to 2nd''?
\end{enumerate}

\problembreak

%-----------------------------------------------------------------------------
\subsection{[75] Hazard Detection Code Analysis}
\hypertarget{prob-4.27}{}
\noindent$\langle$\S4.8$\rangle$\medskip

\noindent Problems in this exercise refer to the following sequence of
instructions, and assume that it is executed on a five-stage pipelined
datapath:
\begin{verbatim}
    add $s3, $s1, $s4
    lw  $s2, 4($s5)
    or  $s2, $s4, $s1
    sw  $s2, 0($s5)
    add $s4, $s2, $s3
\end{verbatim}

\begin{enumerate}
  \phantomsection\addcontentsline{toc}{subsubsection}{4.27.1~[5]}
  \item[4.27.1] [5] If the processor has forwarding, but we forget to
        implement the hazard detection unit, what happens when the original
        code is executed?
  \phantomsection\addcontentsline{toc}{subsubsection}{4.27.2~[10]}
  \item[4.27.2] [10] Change and/or rearrange the code to minimize the number
        of \texttt{nop}(s) needed. You can assume register \texttt{\$s6} can
        be used to hold temporary values in your modified code.
  \phantomsection\addcontentsline{toc}{subsubsection}{4.27.3~[10]}
  \item[4.27.3] [10] If there is forwarding, for the first seven cycles
        during the execution of this code, specify which signals are asserted
        in each cycle by hazard detection and forwarding units.
  \phantomsection\addcontentsline{toc}{subsubsection}{4.27.4~[20]}
  \item[4.27.4] [20] If there is no forwarding, what new input and output
        signals do we need for the hazard detection unit? Using this
        instruction sequence as an example, explain why.
  \phantomsection\addcontentsline{toc}{subsubsection}{4.27.5~[10]}
  \item[4.27.5] [10] For the new hazard detection unit from 4.27.4, roughly
        how many stall cycles are needed for the first 5 cycles during the
        execution of this code?
  \phantomsection\addcontentsline{toc}{subsubsection}{4.27.6~[20]}
  \item[4.27.6] [20] For the new hazard detection unit, what is the total
        number of stall cycles for the entire code sequence?
\end{enumerate}

\problembreak

%-----------------------------------------------------------------------------
\subsection{[60] Branch Prediction CPI Impact}
\hypertarget{prob-4.28}{}
\noindent$\langle$\S4.9$\rangle$\medskip

\noindent Assume the following branch predictor accuracy depends on how often
conditional branches are executed. The breakdown of dynamic instructions into
categories is:

\begin{center}
\begin{tabular}{ccccc}
\toprule
\textbf{R-type} & \textbf{beq/bne} & \textbf{jal} & \textbf{lw} &
\textbf{sw} \\
\midrule
40\% & 25\% & 5\% & 20\% & 10\% \\
\bottomrule
\end{tabular}
\end{center}

\noindent Also, assume the following branch prediction accuracy:

\begin{center}
\begin{tabular}{ccc}
\toprule
\textbf{Always-Taken} & \textbf{Always-Not-Taken} & \textbf{2-Bit} \\
\midrule
45\% & 55\% & 85\% \\
\bottomrule
\end{tabular}
\end{center}

\begin{enumerate}
  \phantomsection\addcontentsline{toc}{subsubsection}{4.28.1~[10]}
  \item[4.28.1] [10] What is the extra CPI due to mispredicted branches with
        the always-taken predictor?
  \phantomsection\addcontentsline{toc}{subsubsection}{4.28.2~[10]}
  \item[4.28.2] [10] Repeat 4.28.1 for the ``always-not-taken'' predictor.
  \phantomsection\addcontentsline{toc}{subsubsection}{4.28.3~[10]}
  \item[4.28.3] [10] Repeat 4.28.1 for the 2-bit predictor.
  \phantomsection\addcontentsline{toc}{subsubsection}{4.28.4~[10]}
  \item[4.28.4] [10] With the 2-bit predictor, what speedup would be
        achieved if we could convert all branch instructions to ALU
        instructions?
  \phantomsection\addcontentsline{toc}{subsubsection}{4.28.5~[10]}
  \item[4.28.5] [10] With the 2-bit predictor, what speedup would be achieved
        if we could convert half of the branch instructions to non-ALU
        instructions?
  \phantomsection\addcontentsline{toc}{subsubsection}{4.28.6~[10]}
  \item[4.28.6] [10] If we know 80\% of all executed branch instructions are
        easy-to-predict loop branches that are always predicted correctly,
        what is the accuracy of the 2-bit predictor on the remaining 20\% of
        branch instructions?
\end{enumerate}

\problembreak

%-----------------------------------------------------------------------------
\subsection{[80] Branch Predictor Design}
\hypertarget{prob-4.29}{}
\noindent$\langle$\S4.9$\rangle$\medskip

\noindent For the following repeating pattern (i.e., in a loop) of branch
outcomes: T, T, NT, T, T, NT, \ldots

\begin{enumerate}
  \phantomsection\addcontentsline{toc}{subsubsection}{4.29.1~[5]}
  \item[4.29.1] [5] What is the accuracy of always-taken and
        always-not-taken predictors for the sequence of branch outcomes?
  \phantomsection\addcontentsline{toc}{subsubsection}{4.29.2~[5]}
  \item[4.29.2] [5] What is the accuracy of the 2-bit predictor for the
        first four branches in this pattern, assuming the predictor starts
        off in the bottom left state from Figure 4.63 (predict not taken)?
  \phantomsection\addcontentsline{toc}{subsubsection}{4.29.3~[10]}
  \item[4.29.3] [10] What is the accuracy of the 2-bit predictor if this
        pattern is repeated forever?
  \phantomsection\addcontentsline{toc}{subsubsection}{4.29.4~[30]}
  \item[4.29.4] [30] Design a predictor that would achieve perfect accuracy
        if this pattern is repeated forever. Your predictor should be a
        sequential circuit with one output that provides a prediction (1 for
        taken, 0 for not taken) and no inputs other than the clock and the
        control signal that indicates the instruction is a conditional branch.
  \phantomsection\addcontentsline{toc}{subsubsection}{4.29.5~[10]}
  \item[4.29.5] [10] What is the accuracy of your predictor from 4.29.4 if
        it is given a repeating pattern that is the exact opposite?
  \phantomsection\addcontentsline{toc}{subsubsection}{4.29.6~[20]}
  \item[4.29.6] [20] Repeat 4.29.4, but now your predictor should be able to
        eventually start perfectly predicting both this pattern and its
        opposite. Your predictor should have an input that tells it the real
        outcome was. Hint: this input lets your predictor determine which of
        the two repeating patterns it is given.
\end{enumerate}

\problembreak

%-----------------------------------------------------------------------------
\subsection{[60] Exception Handling in Pipeline}
\hypertarget{prob-4.30}{}
\noindent$\langle$\S4.10$\rangle$\medskip

\noindent This exercise explores how exception handling affects pipeline
design. The first three problems refer to the following two instructions:

\begin{center}
\begin{tabular}{cc}
\textbf{Instruction 1} & \textbf{Instruction 2} \\
\texttt{beq \$s1, \$s0, label} & \texttt{add \$s3, \$s1, \$s2}
\end{tabular}
\end{center}

\begin{enumerate}
  \phantomsection\addcontentsline{toc}{subsubsection}{4.30.1~[5]}
  \item[4.30.1] [5] Which exceptions can each of these instructions trigger?
        For each exception, specify the pipeline stage in which it is
        detected.
  \phantomsection\addcontentsline{toc}{subsubsection}{4.30.2~[10]}
  \item[4.30.2] [10] If there is a separate handler address for each
        exception, show how the pipeline organization must be changed to
        handle this exception.
  \phantomsection\addcontentsline{toc}{subsubsection}{4.30.3~[10]}
  \item[4.30.3] [10] If the second instruction is fetched immediately after
        the first, describe what happens in the pipeline when the first
        instruction causes an exception. Show the pipeline execution diagram.
  \phantomsection\addcontentsline{toc}{subsubsection}{4.30.4~[20]}
  \item[4.30.4] [20] In vectored exception handling, the table of exception
        handler addresses is in data memory at a known address. Change the
        pipeline to implement this. Repeat Exercise 4.30.3 using this
        modified pipeline.
  \phantomsection\addcontentsline{toc}{subsubsection}{4.30.5~[15]}
  \item[4.30.5] [15] We want to emulate vectored exception handling on a
        machine that has only one fixed handler address. Write the code that
        should be at that fixed address.
\end{enumerate}

\problembreak

%-----------------------------------------------------------------------------
\subsection{[180] Two-Issue Processor}
\hypertarget{prob-4.31}{}
\noindent$\langle$\S4.11$\rangle$\medskip

\noindent In this exercise we compare the performance of 1-issue and 2-issue
processors. Problems refer to the following loop:
\begin{verbatim}
          li   $s1, 0
          j    ENT
    TOP:  add  $s1, $s2, $s0
          lw   $s2, 0($s3)
          sub  $s3, $s3, $s4
    ENT:  bne  $s0, $s1, TOP
\end{verbatim}

\noindent Assume the two-issue, statically scheduled processor has these
properties:
\begin{enumerate}
  \item[1.] One instruction must be a memory operation; the other must be an
        arithmetic/logic instruction or a branch.
  \item[2.] The processor has all possible forwarding paths between stages.
  \item[3.] The processor has perfect branch prediction.
  \item[4.] Two instructions may not issue together in a packet if one
        depends on the other.
  \item[5.] If a stall is necessary, both instructions in the issue packet
        must stall.
\end{enumerate}

\begin{enumerate}
  \phantomsection\addcontentsline{toc}{subsubsection}{4.31.1~[30]}
  \item[4.31.1] [30] Draw a pipeline diagram showing how the MIPS code above
        executes on the two-issue processor.
  \phantomsection\addcontentsline{toc}{subsubsection}{4.31.2~[10]}
  \item[4.31.2] [10] What is the speedup of going from a one-issue to a
        two-issue processor? (Assume the loop runs thousands of iterations.)
  \phantomsection\addcontentsline{toc}{subsubsection}{4.31.3~[10]}
  \item[4.31.3] [10] Rearrange/rewrite the MIPS code to achieve better
        performance on the two-issue processor.
  \phantomsection\addcontentsline{toc}{subsubsection}{4.31.4~[20]}
  \item[4.31.4] [20] Repeat Exercise 4.31.1 using your optimized code from
        4.31.3.
  \phantomsection\addcontentsline{toc}{subsubsection}{4.31.5~[30]}
  \item[4.31.5] [30] Repeat Exercise 4.31.1, but this time use your
        optimized code from Exercise 4.31.3.
  \phantomsection\addcontentsline{toc}{subsubsection}{4.31.6~[10]}
  \item[4.31.6] [10] What is the speedup of going from a one-issue to a
        two-issue processor when running the optimized code?
  \phantomsection\addcontentsline{toc}{subsubsection}{4.31.7~[10]}
  \item[4.31.7] [10] Unroll the MIPS code from Exercise 4.31.3 so that each
        iteration of the unrolled loop handles two iterations of the original
        loop. Then, rearrange/rewrite your unrolled code to achieve better
        performance on the one-issue processor.
  \phantomsection\addcontentsline{toc}{subsubsection}{4.31.8~[20]}
  \item[4.31.8] [20] Repeat Exercise 4.31.4 for the unrolled code on a
        two-issue processor.
  \phantomsection\addcontentsline{toc}{subsubsection}{4.31.9~[10]}
  \item[4.31.9] [10] What is the speedup of going from a one-issue to a
        two-issue processor? (Assume the loop runs thousands of iterations.)
  \phantomsection\addcontentsline{toc}{subsubsection}{4.31.10~[30]}
  \item[4.31.10] [30] What is the speedup if the two-issue processor can run
        two arithmetic/logic instructions together, but the instructions in
        a packet can be any type except that two memory operations cannot be
        scheduled at the same time?
\end{enumerate}

\problembreak

%-----------------------------------------------------------------------------
\subsection{[55] Energy Efficiency Analysis}
\hypertarget{prob-4.32}{}
\noindent$\langle$\S4.3, 4.7, 4.15$\rangle$\medskip

\noindent This exercise explores energy efficiency and its relationship with
performance. Assume the following energy consumption for activity in
Instruction memory, Registers, and Data memory:

\begin{center}
\begin{tabular}{cccccc}
\toprule
\textbf{I-Mem} & \textbf{Reg Read} & \textbf{Reg Write} & \textbf{ALU} &
\textbf{D-Mem Read} & \textbf{D-Mem Write} \\
\midrule
140pJ & 60pJ & 80pJ & 95pJ & 80pJ & 55pJ \\
\bottomrule
\end{tabular}
\end{center}

\begin{enumerate}
  \phantomsection\addcontentsline{toc}{subsubsection}{4.32.1~[5]}
  \item[4.32.1] [5] How much energy is spent to execute an \texttt{add}
        instruction in a single-cycle design and in the five-stage pipelined
        design?
  \phantomsection\addcontentsline{toc}{subsubsection}{4.32.2~[10]}
  \item[4.32.2] [10] What is the worst-case MIPS instruction in terms of
        energy consumption? What is the energy spent to execute it?
  \phantomsection\addcontentsline{toc}{subsubsection}{4.32.3~[10]}
  \item[4.32.3] [10] If energy reduction is paramount, would you change the
        pipeline design? What is the percentage reduction in the energy spent
        by an \texttt{lw} instruction after this change?
  \phantomsection\addcontentsline{toc}{subsubsection}{4.32.4~[10]}
  \item[4.32.4] [10] What other instructions can potentially benefit from
        the change discussed in Exercise 4.32.3?
  \phantomsection\addcontentsline{toc}{subsubsection}{4.32.5~[10]}
  \item[4.32.5] [10] How do your changes from Exercise 4.32.3 affect the
        performance of a pipelined CPU?
  \phantomsection\addcontentsline{toc}{subsubsection}{4.32.6~[10]}
  \item[4.32.6] [10] Can you eliminate the MemRead control signal and have
        the data memory be readable every cycle? Explain why the processor
        will function correctly after this change. If 25\% of instructions
        are loads, what is the effect of this change on clock frequency and
        energy consumption?
\end{enumerate}

\problembreak

%-----------------------------------------------------------------------------
\subsection{[40] Stuck-at Fault Testing}
\hypertarget{prob-4.33}{}
\noindent$\langle$\S4.3, 4.6$\rangle$\medskip

\noindent When silicon chips are fabricated, defects in materials and
manufacturing errors can result in defective circuits. A special class of
cross-talk faults is when a signal is connected to a wire that has a constant
logical value (e.g., power supply wire). These ``stuck-at'' faults cause the
affected signal to always have a logical value of 0 or 1. The following
problems refer to bit 0 of the Write Register input of the Register File.

\begin{enumerate}
  \phantomsection\addcontentsline{toc}{subsubsection}{4.33.1~[10]}
  \item[4.33.1] [10] Assume processor testing is done by running test
        programs and examining results. Can you design a test (values for PC,
        memories, and registers) that would determine if there is a
        stuck-at-0 fault on this signal?
  \phantomsection\addcontentsline{toc}{subsubsection}{4.33.2~[10]}
  \item[4.33.2] [10] Repeat Exercise 4.33.1 for a stuck-at-1 fault.
  \phantomsection\addcontentsline{toc}{subsubsection}{4.33.3~[10]}
  \item[4.33.3] [10] Repeat Exercise 4.33.1, but now the fault is whether
        the \texttt{RegWrite} control signal is always 0 (no fault
        otherwise).
  \phantomsection\addcontentsline{toc}{subsubsection}{4.33.4~[10]}
  \item[4.33.4] [10] Repeat Exercise 4.33.1, but now the fault is whether
        the \texttt{RegWrite} control signal is always 1 (no fault
        otherwise).
\end{enumerate}
